% !TEX root = ../main.tex
\chapter{関数外延性：関数が同じとは何か}
\label{chap:function-extensionality}
\BookIndex{かんすうがいえんせい}{関数外延性}
\BookIndex{funext}{\texttt{funext}}

\begin{chapterabstract}
本章では、関数の等しさを扱います。
ソフトウェアでは、関数、wrapper、adapter を別の実装に置き換える場面が多くあります。
そのときに確認したいのは、名前や内部の書き方が同じかどうかではなく、外から観測できる振る舞いが同じかどうかです。
Lean では、この考え方を「すべての入力に対して同じ出力を返す」という命題として書き、\texttt{funext} を使って関数そのものの等式へ持ち上げます。
本章では、関数外延性を API client wrapper の置き換え安全性として読み替え、実務で使える小さな証明パターンとして整理します。
\end{chapterabstract}

\begin{learninggoals}
\item 関数の等しさを「実装の一致」ではなく「全入力での出力一致」として説明できます。
\item Lean で点ごとの等式 \(\forall x, f\ x = g\ x\) と関数そのものの等式 \(f = g\) の違いを読めます。
\item \texttt{funext} を使って、wrapper や adapter の振る舞い保存を小さく証明できます。
\item 関数外延性で証明できることと、性能、ログ、外部副作用など証明モデル外に残ることを区別できます。
\end{learninggoals}

\section{実装が異なる関数を比較する}

リファクタリングでは、同じ処理を別の書き方に変えることがあります。
たとえば、古い API client を新しい wrapper 経由で呼び出す変更、ロギング用 wrapper の追加、テストダブルによる本番実装の置き換えなどです。

このとき、ソースコードの文字列は変わります。
関数名や、内部で使う補助関数も変わるかもしれません。
しかし、呼び出し側にとって重要なのは、多くの場合、「同じ入力を渡したら同じレスポンスが返るか」です。

図\ref{fig:ch21-overview}は、二つの関数へ同じ任意入力を渡し、出力が点ごとに一致することから関数の等式へ進む流れを示します。

\FigurePlaceholder{fig-ch21-overview.png}{0.90}{関数外延性の全体像}{fig:ch21-overview}

本章は、前章までの道具を次のようにつなぎます。

\begin{itemize}
\item \ChapterRef{chap:ch18_equations_rewriting_normalization}の等式推論では、式や計算結果の一致を扱いました。
\item \ChapterRef{chap:cases-patterns}の場合分けでは、入力の形ごとに証明を分けました。
\item \ChapterRef{chap:induction}の帰納法では、任意の長さのリストや再帰構造を扱いました。
\item 本章では、任意の入力に対する一致を使って、関数そのものの一致を扱います。
\end{itemize}

つまり本章では、「ある入力では同じでした」という主張ではなく、「すべての入力で同じであれば、関数として同じとみなせます」という考え方を扱います。

\section{ソフトウェア上の問題：実装は違うが振る舞いは同じか}

次のような変更を考えます。

\begin{itemize}
\item API client の内部実装を整理しました。
\item wrapper を追加しましたが、レスポンス値は変えない予定です。
\item 旧関数 \texttt{oldClient} を新関数 \texttt{newClient} に置き換えます。
\item テストダブルが、本番実装と同じ抽象仕様を満たしているか確認します。
\end{itemize}

どの場面でも、確認したい性質は共通しています。

\begin{quote}
任意のリクエスト \texttt{r} について、\texttt{oldClient r} と \texttt{newClient r} が同じレスポンスを返す。
\end{quote}

これは、具体的な一件のテストではなく、すべての入力についての主張です。

ただし、注意すべき点があります。
Lean 上の関数が純粋な関数としてモデル化されているなら、同じ入力に対して同じ値を返すことには強い意味があります。
一方、実システムの関数がログ出力、メトリクス送信、時刻参照、乱数、ネットワーク呼び出し、キャッシュ更新などを含む場合、何を「同じ振る舞い」と見るかを先に決める必要があります。

\section{関数外延性の定義}

関数 \(f\) と \(g\) が同じであることを、次のように考えます。

型 \(\alpha\) から型 \(\beta\) への二つの関数 \(f, g : \alpha \to \beta\) について、すべての入力 \(x : \alpha\) で \(f(x) = g(x)\) が成り立つとき、\(f\) と \(g\) は外延的に等しいといいます。
Lean では、点ごとの一致を
\[
  \forall x : \alpha,\ f\ x = g\ x
\]
として表し、関数そのものの等式を
\[
  f = g
\]
として表します。

ここで「外延的」とは、外から観測できる入出力関係に注目する、という意味です。
関数の定義本文が同じかどうか、同じ補助関数を使っているかどうか、同じアルゴリズムかどうかは、外延的等しさの直接の条件ではありません。

関数外延性は、「計算コストも同じです」「ログも同じです」「内部状態の変化も同じです」という意味ではありません。
本章で扱う等しさは、Lean のモデル上で明示した入力と出力の等しさです。
性能や外部副作用まで仕様に含める場合は、それらを戻り値や状態としてモデル化する必要があります。

\section{Lean 4 で点ごとの一致を書く}

まず、関数がすべての入力で同じ結果を返す、という性質を Lean で書きます。

Lean では \texttt{∀ x, f x = g x} が点ごとの等式です。
\texttt{rw} などで関数そのものを置き換えるために \texttt{f = g} が必要なら、\texttt{funext} で橋渡しします。

\begin{listing}[htbp]
\begin{minted}{lean4}
namespace Chapter21

-- 点ごとの一致を、名前を付けて読みやすくする。
def pointwiseEq {α β : Type} (f g : α → β) : Prop :=
  ∀ x : α, f x = g x

-- すべての入力で同じなら、関数そのものも等しい。
theorem funext_from_pointwise
    {α β : Type} {f g : α → β}
    (h : pointwiseEq f g) : f = g := by
  funext x
  exact h x

end Chapter21
\end{minted}
\caption{\texttt{funext\_from\_pointwise}}
\label{lst:ch21-pointwise-funext}
\end{listing}

この証明は、次のように読みます。

\begin{enumerate}
\item 証明したいゴールは \texttt{f = g} です。
\item \texttt{funext x} により、「任意の \texttt{x} について \texttt{f x = g x} を示せばよい」という形に変わります。
\item 仮定 \texttt{h} は、任意の \texttt{x} に対する一致を与えます。
\item したがって、\texttt{exact h x} で証明が終わります。
\end{enumerate}

図\ref{fig:ch21-funext-proof}では、\texttt{funext x} によって関数等式のゴールを、任意の \texttt{x} における出力の等式へ変換します。

\FigurePlaceholder{fig-ch21-funext-proof.png}{0.88}{\texttt{funext} の証明形}{fig:ch21-funext-proof}

\section{実装が違っても、同じ関数になる例}

最初に、非常に小さい例を見ます。
次の二つの関数は、書き方が異なりますが、どちらも自然数に 1 を足します。

\begin{listing}[htbp]
\begin{minted}{lean4}
namespace Chapter21

def oldPlusOne (n : Nat) : Nat :=
  (fun x => x + 1) n

def newPlusOne (n : Nat) : Nat :=
  n + 1

-- 関数そのものが等しいことを示す。
theorem oldPlusOne_eq_newPlusOne :
    oldPlusOne = newPlusOne := by
  funext n
  rfl

end Chapter21
\end{minted}
\caption{\texttt{oldPlusOne\_eq\_newPlusOne}}
\label{lst:ch21-old-new-plus-one}
\end{listing}

\texttt{rfl} は、定義を展開すると同じ形になる場合に使えます。
ここでは、任意の \texttt{n} を取ると、\texttt{oldPlusOne n} も \texttt{newPlusOne n} も、定義展開によって \texttt{n + 1} になります。

この例では、両定義を簡約すると同じラムダ式になるため、直接 \texttt{rfl} でも証明できます。
ここでは、定義上は一致しない関数にも使える一般的な手順を示すため、\texttt{funext n} で点ごとの等式へ分解しています。

\section{no-op wrapper を証明する}

次に、wrapper の例を見ます。
ここでの wrapper は、本来であればログ、メトリクス、ヘッダ追加などを担当する層を表します。
ただし本章では、まず戻り値を変えない wrapper だけをモデル化します。

\begin{listing}[htbp]
\begin{minted}{lean4}
namespace Chapter21

-- 入力をそのまま元の関数に渡すだけの wrapper。
def noopWrapper {α β : Type} (f : α → β) : α → β :=
  fun x => f x

-- wrapper を通しても、関数としては同じ。
theorem noopWrapper_eq {α β : Type} (f : α → β) :
    noopWrapper f = f := by
  funext x
  rfl

end Chapter21
\end{minted}
\caption{\texttt{noopWrapper\_eq}}
\label{lst:ch21-noop-wrapper}
\end{listing}

この定理は、実務では「wrapper を通しても戻り値は変わりません」と読み替えられます。

現実の wrapper がログを書くなら、「ログも観測対象に含めるか」を決める必要があります。
戻り値だけを仕様化する場合、この証明が保証するのは戻り値の保存だけです。
ログ出力まで同じであることは保証しません。

\section{関数そのものの等式と点ごとの等式を使い分ける}

Lean では、次の二つは似ていますが、使う場面が少し異なります。

\begin{itemize}
\item 点ごとの等式：\texttt{∀ x, f x = g x}
\item 関数そのものの等式：\texttt{f = g}
\end{itemize}

点ごとの等式は、ある入力を指定して結果を比較するときに使いやすい形です。
一方、関数そのものの等式は、\texttt{rw} で関数をまとめて置き換えるときに使いやすい形です。

\begin{listing}[htbp]
\begin{minted}{lean4}
namespace Chapter21

theorem apply_pointwise
    {α β : Type} {f g : α → β}
    (h : ∀ x : α, f x = g x) (x : α) :
    f x = g x := by
  exact h x

theorem function_eq_to_pointwise
    {α β : Type} {f g : α → β}
    (h : f = g) :
    ∀ x : α, f x = g x := by
  intro x
  rw [h]

end Chapter21
\end{minted}
\caption{\texttt{apply\_pointwise} と \texttt{function\_eq\_to\_pointwise}}
\label{lst:ch21-pointwise-vs-function}
\end{listing}

この二つを使い分けると、証明が読みやすくなります。
まず点ごとの仕様として「任意の入力で結果が同じです」と書き、必要になったら \texttt{funext} で関数等式にする、という流れが基本です。

\section{ミニケース：新旧 API client wrapper が同じ振る舞いをすること}

ここから、少し実務に近いモデルを作ります。
API client は、リクエストを受け取り、レスポンスを返す関数として表します。

\begin{listing}[htbp]
\begin{minted}{lean4}
namespace Chapter21

structure Request where
  userId : Nat
  amount : Nat
deriving Repr

structure Response where
  status : Nat
  value : Nat
deriving Repr

-- 旧 client。金額に 1 を足して返すだけの単純モデル。
def clientV1 (r : Request) : Response :=
  { status := 200, value := r.amount + 1 }

-- 新 client。補助的な let を使うが、返す値は同じ。
def clientV2 (r : Request) : Response :=
  let calculated := r.amount + 1
  { status := 200, value := calculated }

#eval clientV1 { userId := 1, amount := 100 }  -- => { status := 200, value := 101 }

end Chapter21
\end{minted}
\caption{\texttt{clientV1} と \texttt{clientV2}}
\label{lst:ch21-api-model}
\end{listing}

この例では、\texttt{clientV1} と \texttt{clientV2} の内部実装が少し異なります。
しかし、任意のリクエストに対して返すレスポンスは同じです。

\begin{listing}[htbp]
\begin{minted}{lean4}
namespace Chapter21

-- すべての Request で同じ Response を返すので、
-- 関数としても等しい。
theorem clientV1_eq_clientV2 :
    clientV1 = clientV2 := by
  funext r
  rfl

end Chapter21
\end{minted}
\caption{\texttt{clientV1\_eq\_clientV2}}
\label{lst:ch21-api-client-eq}
\end{listing}

この定理が述べていることは、次の一文に相当します。

\begin{quote}
任意の \texttt{Request} について、\texttt{clientV1} と \texttt{clientV2} は同じ \texttt{Response} を返す。
\end{quote}

つまり、Lean 上でモデル化した範囲では、\texttt{clientV1} から \texttt{clientV2} への置き換えは、戻り値の意味を変えません。

図\ref{fig:ch21-api-wrapper}は、任意の \texttt{Request} を旧 client と新 client へ入力し、同じ \texttt{Response} に到達することを表します。

\FigurePlaceholder{fig-ch21-api-wrapper.png}{0.92}{新旧 API client}{fig:ch21-api-wrapper}

\section{wrapper の順序を変えても戻り値が変わらない例}

次に、二つの wrapper を考えます。
ここでは \texttt{addDefaultHeader} と \texttt{recordMetrics} という名前を付けますが、簡単化のため、どちらも戻り値を変えない wrapper としてモデル化します。

\begin{listing}[htbp]
\begin{minted}{lean4}
namespace Chapter21

def addDefaultHeader
    (client : Request → Response) : Request → Response :=
  fun r => client r

def recordMetrics
    (client : Request → Response) : Request → Response :=
  fun r => client r

theorem wrapper_order_eq (client : Request → Response) :
    addDefaultHeader (recordMetrics client)
      = recordMetrics (addDefaultHeader client) := by
  funext r
  rfl

end Chapter21
\end{minted}
\caption{\texttt{wrapper\_order\_eq}}
\label{lst:ch21-wrapper-order}
\end{listing}

この証明は短いものですが、その保証範囲を正しく読む必要があります。
モデル化した範囲、すなわち「戻り値だけ」を見る限り、二つの wrapper の順序を入れ替えても結果は変わりません。

この定理は、実システムでヘッダ追加とメトリクス記録の順序を常に入れ替えてよい、という意味ではありません。
ヘッダの有無をメトリクスが読む、ログのタイムスタンプが変わる、例外処理の順序が変わる、といった観測を仕様に含めるなら、モデルを拡張する必要があります。
本章の定理は、戻り値に関して順序を入れ替えてよい、という限定された主張です。

\section{場合分けを含む関数外延性}

関数の中に \texttt{Option} や \texttt{Sum} の場合分けがあるときも、基本は同じです。
任意の入力を取り、その入力に対する結果の等式を示します。
ただし、結果が分岐を含む場合は、前章までに学んだ \texttt{cases} が必要になることがあります。

次の \texttt{clientWithFallback} は、\texttt{Option Response} を返す client を受け取り、\texttt{some} ならそのまま返し、\texttt{none} なら \texttt{none} を返します。
つまり、戻り値だけを見ると何も変えていません。

\begin{listing}[htbp]
\begin{minted}{lean4}
namespace Chapter21

def clientWithFallback
    (client : Request → Option Response) : Request → Option Response :=
  fun r =>
    match client r with
    | some res => some res
    | none => none

theorem clientWithFallback_eq
    (client : Request → Option Response) :
    clientWithFallback client = client := by
  funext r
  cases h : client r <;> simp [clientWithFallback, h]

end Chapter21
\end{minted}
\caption{\texttt{clientWithFallback\_eq}}
\label{lst:ch21-option-wrapper}
\end{listing}

この証明では、\texttt{funext r} によって、任意のリクエスト \texttt{r} に対する等式へ移ります。
その後、\texttt{client r} の値を名前付きで場合分けし、\texttt{clientWithFallback} を展開して簡約します。

このように、関数外延性は単独で使うだけでなく、\texttt{cases} や \texttt{simp} と組み合わせて使うことがよくあります。

\section{テストと関数外延性の違い}

通常のテストでは、いくつかの入力を選んで結果を比較します。

\begin{itemize}
\item \texttt{amount = 0} の場合
\item \texttt{amount = 100} の場合
\item \texttt{amount = 9999} の場合
\end{itemize}

これは実務上重要です。
境界値、実データに近いケース、外部システムとの連携は、テストで確認する価値があります。

一方、関数外延性を使った証明では、入力を一つずつ列挙しません。
任意の入力 \texttt{r} を取り、その入力について左右が同じであることを示します。
したがって、モデル化した範囲では、すべての入力に対する主張になります。

テストは、実装上の代表例や外部連携を確認できます。
関数外延性の証明は Lean に書いたモデルの全入力を覆いますが、モデルと実装の対応やモデル外の観測までは保証しません。

\section{実務で使うときの設計手順}

関数外延性を実務に持ち込むときは、最初から大きな実装全体を Lean に移す必要はありません。
次の順で小さく切り出します。

\begin{enumerate}
\item 比較したい二つの関数を決めます。
\item 入力型と出力型を小さなモデルとして定義します。
\item 何を観測対象にするか決めます。
戻り値だけか、ログや状態も含めるかを選びます。
\item 点ごとの仕様 \texttt{∀ x, old x = new x} を書きます。
\item 必要なら \texttt{funext} で \texttt{old = new} にします。
\item 証明した範囲と証明していない範囲を設計文書に明記します。
\end{enumerate}

この手順は、API client だけでなく、正規化関数、DTO converter、権限判定関数、価格計算関数、feature flag の wrapper などにも応用できます。

\section{本章のまとめ}

関数の等しさは、実装本文の一致ではなく、すべての入力に対する出力の一致として捉えられます。

Lean では、点ごとの一致を \texttt{∀ x, f x = g x}、関数そのものの等式を \texttt{f = g} として表します。

\texttt{funext} は、関数そのものの等式を、任意の入力での等式に分解する道具です。

wrapper や adapter の置き換え安全性は、モデル化した範囲では関数外延性の証明として表せます。

証明できるのは、型と関数として明示した観測対象についてだけです。
ログ、性能、外部副作用などを含める場合は、それらもモデルに入れる必要があります。
